\section{An ancient inverse for the flat linearization}
\label{sec:flat}

Fix the profile supplied by Theorem~\ref{prof:profile}, and put
\[
 \rho=\sqrt{-t},\qquad x=|z|^2,\qquad
 h(y)=\frac{G(y)-1}{y},\qquad
 F(z,t)=\rho^{-1}h(|z|^2/\rho),\qquad z\in\R^m,
\]
where $m=24$ and $h(0)=G'(0)$. Throughout this section the time
domain is the entire interval $(-\infty,0)$. The flat linearization is
\begin{equation}\label{flat:operator}
 \widehat{\Pcal}V
 =V_t+\frac{A(z,t)}{16}\Delta_z^2V+B(z,t)V,
 \quad
 A=G(|z|^2/\rho)^2,\quad
 B=2|z|^2G(|z|^2/\rho)\widehat{\A}F,
 \qquad \widehat{\A}=\frac1{16}\Delta_z^2.
\end{equation}
The spaces below are the radial versions of Definition~\ref{ana:spaces},
with scale $\sigma=(|z|^2+\rho)^{1/2}$.

\begin{theorem}[Ancient flat right inverse]\label{flat:inverse}
There is a number $\delta\in(0,1)$ such that, for every fixed
$\alpha\in(0,1)$, the operator \eqref{flat:operator} has a bounded real
linear right inverse
\[
 \widehat D:Y\longrightarrow X,\qquad
 \widehat{\Pcal}\widehat D f=f,\qquad
 \|\widehat D f\|_X\le C\|f\|_Y,
 \qquad p_0=-2+2\delta.
\]
Here $Y$ has weight $p_0-4$ and parabolic exponent $\alpha$, and $X$
has weight $p_0$ and includes the parabolic $\alpha$ norms of all
spatial derivatives through order four and of the first time
derivative. The solution is radial and classical on $t<0$.
Its initial values are selected by the construction. The constant
depends only on the fixed profile, $\delta$, $\alpha$, and the fixed
choices in the local norms.
\end{theorem}

\subsection{Similarity coordinates and coefficient bounds}

The flat equation for $F$ in Theorem~\ref{prof:profile} and the
radial operator identity of Lemma~\ref{geo:operators} yield
\begin{equation}\label{flat:profile-identities}
 \frac1{16}\Delta_\zeta^2\bigl(h(|\zeta|^2)\bigr)
 =-\frac{G'(y)}{2G(y)^2},\qquad y=|\zeta|^2,
 \qquad F_t+G(|z|^2/\rho)^2\widehat{\A}F=0.
\end{equation}
In particular, with
\[
 A_f(\zeta)=G(|\zeta|^2)^2,\qquad
 B_f(\zeta)
 =\frac{2|\zeta|^2G(|\zeta|^2)}{16}
       \Delta_\zeta^2\bigl(h(|\zeta|^2)\bigr)
 =-\frac{yG'(y)}{G(y)},
\]
one has $A(z,t)=A_f(z/\sqrt\rho)$ and
$B(z,t)=\rho^{-2}B_f(z/\sqrt\rho)$.
These coefficients are smooth at the origin. Positivity and the
differentiated asymptotic expansion in Theorem~\ref{prof:profile} imply
\begin{equation}\label{flat:symbols}
 \inf A_f>0,\qquad
 |\partial_\zeta^\nu(A_f-c^2)|+
 |\partial_\zeta^\nu B_f|
 \le C_\nu(1+|\zeta|)^{-4-|\nu|}.
\end{equation}
In particular all derivatives, including any fixed number of
applications of $\zeta\cdot\nabla$, are bounded.

Set
\[
 \tau=-\log(-t),\qquad \zeta=z/\sqrt\rho,
 \qquad V(z,t)=\rho^{-1}v(\zeta,\tau),
 \qquad f^\#(\zeta,\tau)=\rho^3f(\sqrt\rho\,\zeta,t).
\]
Since $\dd\tau/\dd t=\rho^{-2}$ and
$\partial_t\zeta=\zeta/(4\rho^2)$, direct differentiation gives
\[
 V_t=\rho^{-3}
   \left(v_\tau+\tfrac14\zeta\cdot\nabla v+\tfrac12v\right),
 \qquad \Delta_z^2V=\rho^{-3}\Delta_\zeta^2v.
\]
Thus \eqref{flat:operator} becomes the autonomous equation
\begin{equation}\label{flat:similarity-equation}
 v_\tau+\frac{A_f}{16}\Delta^2v
       +\frac14\zeta\cdot\nabla v+(\tfrac12+B_f)v=f^\#.
\end{equation}
We shall first solve it in
$E=C_{0,\mathrm{rad}}(\R^m;\C)$, with the supremum norm. This denotes
continuous radial functions tending to zero at spatial infinity;
complex scalars are needed only for spectral theory.

\subsection{Evolution on continuous functions vanishing at infinity}

We give the evolution construction since the unbounded drift in
\eqref{flat:similarity-equation} and the change in the principal
coefficient both matter. On a time interval beginning at $\tau_0$, put
\[
 u(\eta,s)=v(e^{s/4}\eta,\tau_0+s),\qquad 0\le s\le1.
\]
Then
\begin{equation}\label{flat:moving-equation}
 u_s+a(\eta,s)\Delta_\eta^2u+b(\eta,s)u=l(\eta,s),
 \qquad
 a=\frac{e^{-s}}{16}A_f(e^{s/4}\eta),\quad
 b=\tfrac12+B_f(e^{s/4}\eta).
\end{equation}
The right side is the corresponding dilation of $f^\#$.
By \eqref{flat:symbols}, $a,b$ have bounded derivatives of every order
on this slab, and $0<\kappa\le a\le\kappa^{-1}$ for a fixed
$\kappa$. The following construction also applies without radial
symmetry. It is the scalar whole-space version of the higher-order
heat-kernel parametrix; see also \cite[Appendix~A]{HeZeng2021}.

Let
\[
 k_\theta(w)=(2\pi)^{-m}\int_{\R^m}
               e^{iw\cdot\xi-\theta|\xi|^4}\,\dd\xi,
 \qquad \theta>0.
\]
Its integral is one. Scaling and integration by parts in the Fourier
integral show that, uniformly for $\kappa\le a_*\le\kappa^{-1}$,
\begin{equation}\label{flat:kernel-bounds}
 |\partial_{a_*}^{\ell}\partial_w^\nu k_{a_*q}(w)|
 \le C_{N,\ell,\nu}\,
 q^{-(m+|\nu|)/4}
       (1+|w|/q^{1/4})^{-N},\qquad 0<q\le1.
\end{equation}
Indeed after scaling $\xi=q^{-1/4}\xi'$ the multipliers and all their
derivatives form a bounded set of Schwartz functions. No sign property
of $k_\theta$ is used.

Freeze the principal coefficient at the source point and time:
\[
 K_{s,r}(\eta,y)=k_{a(y,r)(s-r)}(\eta-y),\qquad r<s,
\]
and also denote by $K_{s,r}$ the associated integral operator. Applying
the left side of \eqref{flat:moving-equation} at the evaluation point
gives the residual kernel
\begin{equation}\label{flat:residual}
 J_{s,r}(\eta,y)
 =\bigl(a(\eta,s)-a(y,r)\bigr)
       \Delta_\eta^2 k_{a(y,r)(s-r)}(\eta-y)
       +b(\eta,s)K_{s,r}(\eta,y).
\end{equation}
The global Lipschitz bound for $a$ and \eqref{flat:kernel-bounds} give
\begin{equation}\label{flat:residual-bound}
 \|K_{s,r}\|_{C_0\to C_0}\le C,
 \qquad
 \|J_{s,r}\|_{C_0\to C_0}
 \le C\bigl((s-r)^{-3/4}+1\bigr)
 \le C'(s-r)^{-3/4}.
\end{equation}
For the first residual term, the spatial difference $|\eta-y|$
contributes $(s-r)^{1/4}$ against a fourth derivative of size
$(s-r)^{-1}$; the temporal difference contributes $(s-r)$.
The kernel bounds first show that both operators preserve $C_0$ for
compactly supported continuous inputs, by decay in $\eta$, and then
for all $C_0$ inputs by uniform approximation.

These kernels give a strong approximate identity. To see this
explicitly, replace $a(y,r)$ by $a(\eta,r)$ in $K_{s,r}$.
The operator norm of the difference is $O((s-r)^{1/4})$, by the
parameter derivative bound in \eqref{flat:kernel-bounds}. For the
replacement kernel, its integral in $y$ is one at each $\eta$;
uniform continuity of a $C_0$ function, followed by the uniform kernel
tail bound, proves convergence to that function in supremum norm.
Consequently
\begin{equation}\label{flat:approximate-identity}
 K_{s,r}d\longrightarrow d\quad\hbox{in }C_0
 \quad\hbox{as }s-r\downarrow0,
\end{equation}
uniformly in the initial time $r$. This is strong convergence; only
the difference between the two frozen kernels was estimated in
operator norm.

For $d\in C_0$ and $l\in C([0,1];C_0)$ solve
\begin{equation}\label{flat:volterra}
 \begin{split}
 j(s)+\int_0^sJ_{s,r}j(r)\,\dd r&=l(s)-J_{s,0}d,\\
 u(s)&=K_{s,0}d+\int_0^sK_{s,r}j(r)\,\dd r.
 \end{split}
\end{equation}
All integrals are Bochner integrals. The inverse of the first equation
is its Volterra series. For completeness, if
$H(s)=Cs^{-3/4}$, its $k$-fold convolution on the positive half-line is
\[
 H^{*k}(s)
 =\frac{C^k\Gamma(1/4)^k}{\Gamma(k/4)}s^{k/4-1}
 \quad(k\ge1).
\]
This identity follows by the beta integral and induction. Since the
right side of the first equation in \eqref{flat:volterra} is bounded
by $Cs^{-3/4}(\|d\|_\infty+\|l\|_{C([0,1];C_0)})$, the same formula
and the growth of the gamma function prove convergence of the series
and the estimates
\begin{equation}\label{flat:evolution-bound}
 \|j(s)\|_\infty\le Cs^{-3/4}
       (\|d\|_\infty+\|l\|_{C([0,1];C_0)}),\qquad
 \sup_{0\le s\le1}\|u(s)\|_\infty
 \le C(\|d\|_\infty+\|l\|_{C([0,1];C_0)}).
\end{equation}
The integral in the formula for $u$ tends to zero as $O(s^{1/4})$.
Together with \eqref{flat:approximate-identity}, this proves
$u\in C([0,1];C_0)$ and $u(0)=d$. Continuity at positive times follows
by dominated convergence, splitting off a short interval adjacent to
the diagonal.

The kernel identity \eqref{flat:residual} proves that $u$ solves
\eqref{flat:moving-equation} distributionally: apply the equation to
the second line of \eqref{flat:volterra}, obtaining
$J_{s,0}d+j(s)+\int_0^sJ_{s,r}j(r)\,\dd r=l(s)$.
This calculation can be made first with a positive separation from
the diagonal and smooth compactly supported tests; the diagonal
term converges to $j$ by \eqref{flat:approximate-identity}, and
\eqref{flat:residual-bound} justifies the remaining integrals.
Lemma~\ref{ana:schauder} and interior regularity then give smoothness
for positive times in the homogeneous case, and classical local
regularity for locally parabolic H\"older sources.

We also record uniqueness in the class just constructed. If $u$ is
the difference of two such solutions, it solves the homogeneous
equation with zero $C_0$ initial value. Let
$W(\eta)=\exp(-\sqrt{1+|\eta|^2})$. Its derivatives are bounded by
constant multiples of $W$. For $s\ge\varepsilon>0$, interior
regularity supplies globally bounded spatial derivatives of $u$.
Integration by parts, first with compact spatial cutoffs, gives
\[
 \Re\int W\overline u\,a\Delta^2u
 =\Re\int\Delta(aW\overline u)\Delta u
 \ge\frac\kappa2\|\Delta u\|_{L^2(W)}^2
                  -C\|u\|_{L^2(W)}^2.
\]
Here the mixed terms are absorbed using
\[
 \|\nabla u\|_{L^2(W)}^2
 \le\epsilon\|\Delta u\|_{L^2(W)}^2
                 +C_\epsilon\|u\|_{L^2(W)}^2.
\]
This last inequality follows by integrating
$\int W|\nabla u|^2$ once by parts and using
$|\nabla W|\le CW$ and Young's inequality. Exponential decay of $W$
allows removal of the cutoffs. The bounded zeroth order coefficient
therefore gives
\[
 \frac{\dd}{\dd s}\|u(s)\|_{L^2(W)}^2
 \le C\|u(s)\|_{L^2(W)}^2.
\]
Gronwall's inequality, followed by $\varepsilon\downarrow0$, proves
$u=0$, because the zero $C_0$ initial trace implies a zero weighted
$L^2$ trace. This also proves uniqueness for the inhomogeneous
construction by subtraction.

Undoing the dilation now gives a strongly continuous homogeneous
semigroup $S(s)$ on $C_0(\R^m)$ for
\eqref{flat:similarity-equation}; concatenation is justified by the
uniqueness just proved. Dilation is strongly continuous on $C_0$,
and the equation is autonomous. On $0\le s\le1$ one has
$\|S(s)\|\le C$. Inhomogeneous evolution for a continuous $C_0$ source
is
\begin{equation}\label{flat:duhamel}
 v(\tau_0+s)=S(s)v(\tau_0)
       +\int_0^sS(s-r)f^\#(\tau_0+r)\,\dd r.
\end{equation}
The formula follows either from the distributional equation and
uniqueness or by approximation with smooth sources. Rotation and
complex conjugation preserve the equation and the construction.
In particular $S(s)$ preserves $E$ and its real subspace.

\subsection{Compactness at positive time}

Write $P=S(1)|_E$, reserving this symbol here for a bounded operator,
and let $P_c$ be the corresponding operator with $A_f=c^2$ and
$B_f=0$. We prove
\begin{equation}\label{flat:compact-difference}
 P-P_c:E\longrightarrow E\quad\hbox{is compact}.
\end{equation}
Both time-one evolutions map their unit balls to locally
equicontinuous functions by the interior estimates. It remains to
prove uniform decay of their difference at spatial infinity.

Suppose this decay fails. There are $d_j\in E$ with
$\|d_j\|_\infty\le1$, points $\zeta_j$ with $|\zeta_j|\to\infty$,
and $\epsilon_0>0$ such that
$|(P-P_c)d_j(\zeta_j)|\ge\epsilon_0$.
In the moving coordinates \eqref{flat:moving-equation}, translate
by $\eta_j=e^{-1/4}\zeta_j$. The two solutions, denoted by
$u_j$ and $u_{c,j}$ after translation, have the same initial value
$d_j(\eta+\eta_j)$ and uniform supremum bounds. Their coefficients
converge, smoothly on compact subsets of
$\R^m\times[0,1]$, to
\[
 a_c(s)=\frac{c^2e^{-s}}{16},\qquad b_c=\tfrac12.
\]
This follows from \eqref{flat:symbols}; it also controls time
derivatives in the moving coordinates.

After passing to a subsequence, the translated initial functions
converge weakly star in $L^\infty(\R^m)$ to some $d_\infty$.
Sequential extraction is available because $L^1(\R^m)$ is separable.
Interior estimates and a diagonal extraction give smooth local
limits $u_\infty,u_{c,\infty}$ for $s>0$, including convergence of
their values at $s=1$. Both limits are bounded on the whole slab.
To identify their initial traces, for
$\psi\in C_c^\infty(\R^m)$ integrate the translated equation:
\begin{equation}\label{flat:translated-trace}
 \left|\int u_j(s)\psi
       -\int d_j(\eta+\eta_j)\psi\right|
 \le Cs\,
 \sup_{j,\,0\le r\le1}
 \left(\|\Delta^2(a_j(\cdot,r)\psi)\|_{L^1}
                    +\|b_j(\cdot,r)\psi\|_{L^1}\right)
 \le C_\psi s.
\end{equation}
All four spatial derivatives have fallen on the coefficient times
the test function. Uniform coefficient bounds make the last
constant independent of $j$. The identity extends down to $s=0$
using the individual $C_0$ initial traces. The same estimate holds
for $u_{c,j}$. Passing to the limit proves that both limits have
distributional initial value $d_\infty$.

Here is the required whole-space uniqueness for these bounded
limits; a $C_0$ initial trace of the limits is not needed.
Their difference $w$ solves
\[
 w_s+a_c(s)\Delta^2w+\tfrac12w=0
\]
and has zero distributional initial value. For a fixed $t_1\in(0,1]$
and $\psi\in\mathcal S(\R^m)$, use the backward test function
\[
 \phi(s)=e^{-(t_1-s)/2}
       k_{\int_s^{t_1}a_c(r)\,\dd r}*\psi,
 \qquad 0\le s\le t_1,
\]
with $\phi(t_1)=\psi$. It solves
$-\phi_s+a_c(s)\Delta^2\phi+\tfrac12\phi=0$, is continuous in the
Schwartz topology, and hence $\int w(s)\phi(s)$ is constant for
$s>0$. One may justify the pairing by compact cutoffs and then use
the Schwartz decay. Boundedness of $w$ extends its zero initial
trace from compactly supported tests to every Schwartz test;
$\phi(s)\to\phi(0)$ in $L^1$ also permits this changing test.
The constant pairing is therefore zero. Since $\psi$ is arbitrary,
$w(t_1)=0$. This contradicts
$|u_\infty(0,1)-u_{c,\infty}(0,1)|\ge\epsilon_0$.

We have proved uniform vanishing of the tails of
$(P-P_c)$ on the unit ball. Together with local equicontinuity,
the Arzel\`a--Ascoli criterion for $C_0$ proves
\eqref{flat:compact-difference}: restrict to a large ball, extract
there, and use the common small tails to obtain uniform convergence
on the whole space. The closed radial subspace is preserved.

\subsection{A spectral weight and the two-sided recursion}

For the constant coefficient equation, the moving-coordinate formula
on an interval of length $k$ is explicit:
\begin{equation}\label{flat:constant-powers}
 (P_c^kd)(\zeta)
 =e^{-k/2}(k_{\theta_k}*d)(e^{-k/4}\zeta),
 \qquad
 \theta_k=\frac{c^2}{16}(1-e^{-k}),\qquad k\ge1.
\end{equation}
The $L^1$ norm of $k_\theta$ is independent of $\theta>0$ by
scaling. Thus $\|P_c^k\|\le Ce^{-k/2}$, and the Neumann resolvent
series converges for $|\lambda|>e^{-1/2}$.

We spell out the spectral consequence of
\eqref{flat:compact-difference}. On the connected domain
$\Omega=\{|\lambda|>e^{-1/2}\}$ factor
\[
 \lambda I-P
 =\left[I-(P-P_c)(\lambda I-P_c)^{-1}\right]
                         (\lambda I-P_c).
\]
The operator subtracted from the identity is compact and analytic
in $\lambda$. The analytic Fredholm alternative implies that
$(\lambda I-P)^{-1}$ is meromorphic on $\Omega$, with finite rank
principal parts, since it is invertible for large $|\lambda|$.
One way to obtain this alternative is to approximate the compact
operator at any one parameter by a finite rank operator, invert the
remaining operator of norm less than one in a neighborhood, and
reduce invertibility to a finite-dimensional analytic determinant.
The required approximation is available here: identify radial $E$
with $C_0([0,\infty))$, and use bounded polygonal interpolation
operators with compact support, which converge strongly to the
identity and hence uniformly on compact subsets of $E$.
The determinant either vanishes identically or has isolated zeros;
continuation through overlapping neighborhoods and invertibility
at a large parameter exclude the first alternative. The same
reduction makes every coefficient of a pole finite rank.
It follows that the spectrum in $\Omega$ consists of isolated
eigenvalues of finite algebraic multiplicity, with no accumulation
inside $\Omega$; these are the usual compact-operator and resolvent
arguments of \cite[Chapter~VII, Section~1.3, and the supplementary notes]{Kato1995}.

Choose
\begin{equation}\label{flat:weight-choice}
 e^{-1/2}<q<1,\qquad
 \{|\lambda|=q\}\cap\spec(P)=\varnothing,\qquad
 \gamma=-\log q\in(0,\tfrac12),\qquad \delta=2\gamma.
\end{equation}
Such a circle exists by the discreteness just proved. There are only
finitely many eigenvalues outside it: they lie in a bounded closed
annulus contained in $\Omega$. Let $\Pi_u$ be their summed Riesz
projection, defined by resolvent integration around these
eigenvalues, and put $\Pi_s=I-\Pi_u$. The resolvent identity gives
commuting projections and the invariant decomposition
\[
 E=E_s\oplus E_u,\qquad E_s=\Pi_sE,\quad E_u=\Pi_uE,
 \qquad P_s=P|_{E_s},\quad P_u=P|_{E_u}.
\]
The space $E_u$ is finite dimensional, $P_u$ is invertible, and
the two restricted spectra lie strictly inside and strictly outside
the chosen circle, respectively. If $E_u=\{0\}$, all terms involving
it below are omitted. Compactness of the restricted spectra and the
spectral radius formula give numbers $r_s<q<r_u$ such that
\begin{equation}\label{flat:dichotomy-bounds}
 \|P_s^k\|\le C r_s^k,\qquad
 \|P_u^{-k}\|\le C r_u^{-k},\qquad k\ge0.
\end{equation}
For example choose $r_s$ strictly between $r(P_s)$ and $q$, and
$r_u$ strictly between $q$ and the smallest modulus of an eigenvalue
of $P_u$; the spectral radius formula controls all sufficiently
large powers, and their finitely many predecessors are absorbed
in $C$. The projections commute with complex conjugation because
the enclosed spectral sets are invariant under conjugation.

Fix this $\delta$. If $M=\|f\|_Y$, the source weight gives the exact
estimate
\begin{equation}\label{flat:source-decay}
 |f^\#(\zeta,\tau)|
 \le M\rho^\delta(1+|\zeta|^2)^{-3+\delta}
 =M e^{-\gamma\tau}(1+|\zeta|^2)^{-3+\delta}.
\end{equation}
Indeed $\sigma^2=\rho(1+|\zeta|^2)$ and
$\rho^3\sigma^{p_0-4}=\rho^\delta(1+|\zeta|^2)^{-3+\delta}$.
This proves that $f^\#(\tau)\in E$. It also proves continuity with
values in $E$: on any compact $\tau$ interval its spatial tails are
uniformly small by \eqref{flat:source-decay}, while on a fixed spatial
ball local H\"older continuity of $f$ and smoothness of the change of
coordinates give uniform continuity in $\tau$.

Let the one-step source increments be
\[
 l_j=\int_0^1S(1-s)f^\#(j+s)\,\dd s,\qquad j\in\Z.
\]
Then $\|l_j\|_E\le CMq^j$. Define
\begin{equation}\label{flat:recursion-solution}
 v_j=\sum_{k=0}^\infty P_s^k\Pi_s l_{j-1-k}
       -\sum_{k=0}^\infty P_u^{-k-1}\Pi_u l_{j+k}.
\end{equation}
Both series converge absolutely for every integer $j$. More
explicitly, their norms are at most
\[
 CMq^j\left[
   q^{-1}\sum_{k=0}^\infty(r_s/q)^k
       +r_u^{-1}\sum_{k=0}^\infty(q/r_u)^k\right]
 \le C'Mq^j.
\]
Shifting the summation indices proves
$v_{j+1}=Pv_j+l_j$. On $[j,j+1]$ define $v$ by
\eqref{flat:duhamel} with initial value $v_j$.
The definitions agree at the endpoints and give
\begin{equation}\label{flat:preliminary-bound}
 \|v(\tau)\|_E\le CM e^{-\gamma\tau},\qquad
 |V(z,t)|\le CM\rho^{-1+\delta}.
\end{equation}
The continuous joining creates no distributional time-boundary
term, so \eqref{flat:similarity-equation} holds across every joined
face. Interior regularity makes it a classical equation there.
The construction is linear and real on real sources. The second
sum in \eqref{flat:recursion-solution} selects the outside spectral
components using future forcing; its use is permitted precisely
because no initial values have been imposed.

\subsection{The full weighted parabolic estimate}

We finish by strengthening \eqref{flat:preliminary-bound} to the
$X$ estimate. This step uses $\delta>0$ as well as $\delta<1$.
For a box centered at $(z_0,t_0)$ with
$r_0^2=|z_0|^2\le\rho_0$, one has
$\sigma_0^2\asymp\rho_0$. On a fixed enlargement of this box,
chosen small in units $(\sigma_0,\sigma_0^4)$, the time stays
strictly negative, $\rho\asymp\rho_0$, and
$|z|^2/\rho$ stays bounded. In these units the coefficients
$A$ and $\sigma_0^4B$ have bounded derivatives of every fixed order,
and the principal coefficient has a uniform positive lower bound.
The source has scaled $C^\alpha$ norm at most $CM$, after division
by $\sigma_0^{p_0-4}$. Also
\[
 |V|\le CM\rho_0^{-1+\delta}
          \le C'M\sigma_0^{p_0}.
\]
Lemma~\ref{ana:schauder}, with the fixed preceding-time enlargement,
therefore gives the required full $X$ bound on every such box.

For the remaining boxes a separate estimate restores the spatial
weight. Fix $a>0$ and write
\[
 z=ay,\qquad t=a^4s,\qquad
 U(y,s)=a^{-p_0}V(ay,a^4s),\qquad
 g(y,s)=a^{4-p_0}f(ay,a^4s).
\]
Use the fixed annuli
\[
 \mathcal U=\{\tfrac12<|y|<2\},\qquad
 \mathcal U'=\{\tfrac34<|y|<\tfrac54\}.
\]
On $\mathcal U\times(-4,0)$ the rescaled equation is
\begin{equation}\label{flat:annular-equation}
 U_s+\frac{G(|y|^2/\sqrt{-s})^2}{16}\Delta_y^2U
      +\mathfrak b(y,s)U=g,\qquad
 \mathfrak b(y,s)=(-s)^{-1}
                     B_f\bigl(y/(-s)^{1/4}\bigr).
\end{equation}
All coefficients are independent of $a$. Their smooth bounds remain
uniform as $s\uparrow0$. To verify the assertion for the zeroth order
term explicitly, put $Y=|y|^2/\sqrt{-s}$ and use
\[
 \mathfrak b(y,s)=-(-s)^{-1}\frac{YG'(Y)}{G(Y)}.
\]
The expansion of $G$ in powers of $Y^{-2}$ gives an expansion of
$-YG'/G$ starting at $Y^{-2}$. Since
$Y^{-2}=(-s)/|y|^4$, division by $-s$ leaves a smooth expansion in
nonnegative integral powers of $-s$. Differentiated remainders of
arbitrarily high order show that both this coefficient and the
principal coefficient extend smoothly to $s=0$ on the annulus.
Away from $s=0$ their bounds follow from ordinary smoothness of $G$.
The principal coefficient stays uniformly positive throughout.

Furthermore
\[
 \frac{\sigma(ay,a^4s)}a
       =\bigl(|y|^2+\sqrt{-s}\bigr)^{1/2}
\]
is bounded above and below on the larger annular slab. A fixed
covering by the local boxes of Definition~\ref{ana:spaces} therefore
gives $\|g\|_{C^\alpha(\mathcal U\times(-4,0))}\le CM$.
Here and below the norm is parabolic; pairs outside a common small
box are controlled by the supremum bound and their fixed positive
separation.

The preliminary estimate becomes
\[
 |U(y,s)|\le CM(-s)^{(-1+\delta)/2},
\]
with no factor of $a$: indeed
$a^{-p_0}(a^2\sqrt{-s})^{-1+\delta}
=(-s)^{(-1+\delta)/2}$. Consequently
\begin{equation}\label{flat:annular-l2}
 \|U\|_{L^2(\mathcal U\times(-4,0))}^2
 \le C M^2|\mathcal U|\int_0^4 q^{-1+\delta}\,\dd q
 =C M^2|\mathcal U|\frac{4^\delta}{\delta}<\infty.
\end{equation}
The value of $\delta$ has already been fixed, so the dependence on
$1/\delta$ is harmless.

Apply Lemma~\ref{ana:l2} to \eqref{flat:annular-equation} on nested
annuli, with the preceding-time buffer from $-4$ to $-2$. It gives
\begin{equation}\label{flat:annular-full}
 \sum_{j=0}^4\|D_y^jU\|_{C^\alpha(\mathcal U'\times(-2,0))}
 +\|U_s\|_{C^\alpha(\mathcal U'\times(-2,0))}
 \le CM.
\end{equation}
One first applies the estimate with terminal time $b<0$. On each
such truncated slab $U$ has finite supremum norm by
\eqref{flat:preliminary-bound}; the coefficient, source, and $L^2$
bounds above are independent of $b$. Letting $b\uparrow0$ then
proves \eqref{flat:annular-full}, including H\"older quotients between
times arbitrarily close to zero. The estimate uses preceding times
and requires no terminal value at $s=0$.

Finally take an arbitrary remaining center with $r_0^2>\rho_0$ and
choose $a=r_0$. Its rescaled time satisfies $-1<s_0<0$ and
$1<\sigma_0/a<\sqrt2$. The entire defining $X$ box lies in
$\mathcal U'\times(-2,0)$ when the fixed box radius $a_0$ is small:
its spatial radius is at most $\sqrt2a_0$ and its time radius at
most $4a_0^4$. In particular a box near the dividing time $s=-1$
has the required preceding-time buffer. Rescaling
\eqref{flat:annular-full} gives its weighted $X$ norm at most $CM$.
The core and annular cases cover all centers. Thus
$\|V\|_X\le C\|f\|_Y$, completing the proof of
Theorem~\ref{flat:inverse}.
